%% kapa: exam=1 round=30 year=2019 period=2 form=A subject=appraisal_law subject_ko=감정평가관계법규 seq=20 number=20
%% source: Q-Net 감정평가사 기출문제 (cst00309) artlSeq=5205205 file=extracted/1008/5205205/2019년 제30회 감정평가사 1차시험 문제지/2019년 제30회 감정평가사 1차시험 문제지/2019년도 제30회 감정평가사 1차 2교시 A형.hwp
%% conversion: kordoc 4.24.51 -> tools/gen1.py (auto); review=proofread-sonnet; flags=none
\begin{question}{20}
국유재산법령상 국유재산에 관한 설명으로 옳지 않은 것은?
\begin{choices}
  \item 사권이 설정된 재산을 판결에 따라 취득하는 경우 그 사권이 소멸된 후가 아니면 국유재산으로 취득하지 못한다.
  \item 국유재산의 범위에는 선박, 지상권, 광업권, 특허권, 저작권이 포함된다.
  \item 총괄청은 다음 연도의 국유재산의 관리ㆍ처분에 관한 계획의 작성을 위한 지침을 매년 4월 30일까지 중앙관서의 장에게 통보하여야 한다.
  \item 총괄청은 일반재산을 보존용재산으로 전환하여 관리할 수 있다.
  \item 확정판결에 따라 일반재산에 사권을 설정할 수 있다.
\end{choices}
\end{question}
\answer{1}
